\documentclass[conference]{IEEEtran}
\IEEEoverridecommandlockouts
% The preceding line is only needed to identify funding in the first footnote. If that is unneeded, please comment it out.
\usepackage{cite}
\usepackage{amsmath,amssymb,amsfonts}
\usepackage{algorithmic}
\usepackage{graphicx}
\usepackage{textcomp}
\usepackage{xcolor}
\def\BibTeX{{\rm B\kern-.05em{\sc i\kern-.025em b}\kern-.08em
    T\kern-.1667em\lower.7ex\hbox{E}\kern-.125emX}}
\begin{document}

\title{Textual Temporal Analysis of a Presidential Term
}

\author{\IEEEauthorblockN{1\textsuperscript{st} Roselyn Crisostomo}
\IEEEauthorblockA{\textit{Department of Information Systems and Computer Science} \\
\textit{Ateneo De Manila University}\\
Quezon City, Philippines \\
roselyn.mae.crisostomo@student.ateneo.edu}
}

\maketitle

\begin{abstract}
Abstract
\end{abstract}

\begin{IEEEkeywords}
social computing
\end{IEEEkeywords}

\section{Introduction}
Temporal analysis is a way to gain insight into a data's behavior over time. It allows examiners to see patterns and understand which could influence a certain variable. Most commonly it is used for predicting future outcomes based on past behavior or trends, such as forecasting sales. In this research, it will be applied on a dataset of Philippine legislative statutes or Republic Acts (RA), representing a snapshot of the Former President Rodrigo Roa Duterte's (FPRRD) presidency.

There are multiple facets in appreciating a past leader's performance. It can be judged based on the economic impact or even on mass popularity. This research focuses on identifying the tangible effects under the former president's administration by looking into the new laws or amendments performed during his term. This analysis also seeks to objectively evaluate whether there were effective legal reforms that reflected what was promised to the Filipino people. The study will apply social computing techniques such as topic modeling and clustering analysis on the collected dataset to process the textual data, then provide an interpretation of the outcome of those models.

\section{Research Question}
This study seeks to answer the following research question:

\begin{enumerate}
    \item How did FPRRD implement his objectives stated from each of his State of the Nation Address (SONA) from 2016 to 2021?
\end{enumerate}

\section{Review of Related Literature}

\subsection{Text Analysis in Law}
Legislation is the exercise of authority by the government to address issues or improvements of the nation. The number of legislations naturally accumulate over time, and these voluminous documents reflect what state of the country at a certain period. These also give pertinent value when analyzing history through the lens of lawmaking.

Modern technology has significantly assisted with working on numerous text documents and one example is by utilizing graph neural networks and mapping relationships among various documents. San Pedro and Ballera implemented this for the purpose of constructing a legal knowledge graph of Philippine RAs. It was a statutory analysis of the most cited statutes among 12,124 RAs, and confirmed that RA 7160, the Local Government Code of 1991 is the most influential statute.

One kind of application for text analysis in the legal field is argument mining. It applies Natural Language Processing (NLP) to identify structures and categorize as premise, conclusion, and such. In legal documents, Poudyal, Gonçalves, and Quaresma developed a clustering technique for implementing it which aims to provide a solution to an issue where one argument may be part of multiple other distinct arguments.

Lastly, Malandrino presented a comparative study between qualitative and quantitive studies for Social Network Analysis (SNA), specifically for policy networks (PN). SNA aims to reveal observable relationships among different entities or actors, and the comparative study was used to understand how PNs are interacted with by those actors, such as their coordination. Applying NLP quantitatively (i.e., finding patterns of frequencies) guides the researcher through general text patterns, which complements the qualitative approach (i.e., topic modeling) which hints at nuanced or deeper meaning behind the numbers.

Overall, there are plenty of different ways to apply text analysis in the context of legal documents which enables researchers to gain deeper understanding on their content and structures.

\subsection{Text Analysis in Presidential Speeches}
SONA is an annual presidential speech delivered by the president of the Philippines. It informs the public about the administration's merits and accomplishments, as well as any future plans.

Presidential speeches are a good way to analyze rhetorics and speech behaviors of leaders, and assess how they sway public opinions with what they include and exclude from their speech, as well as how they say certain things such as choosing certain vocabulary over another. In Ancho, et al.'s analysis of the speeches of President Arroyo, President Aquino, and President Duterte, specifically in the context of education and how they addressed it, they used word frequency and thematic analysis using Nvivo software.

Miranda and Bringula's approach leans more into revealing the underlying patterns from multiple presidential speeches spanning from the past 13 Philippine presidents. They used topic modelling to extract prominent themes as well as term frequency to compare the priorities of each president. They also used sentiment analysis to classify whether each speech was positive or negative by the start and end of each year, and by how much.

These existing studies serve as guides on how to interpret the SONAs of FPRRD which will be focused on in this research. It also acts as reference with what themes to look out for during the analysis.

\section{Methodology}


\begin{thebibliography}{00}
\bibitem{b1} G. Eason, B. Noble, and I. N. Sneddon, ``On certain integrals of Lipschitz-Hankel type involving products of Bessel functions,'' Phil. Trans. Roy. Soc. London, vol. A247, pp. 529--551, April 1955.
\bibitem{b2} J. Clerk Maxwell, A Treatise on Electricity and Magnetism, 3rd ed., vol. 2. Oxford: Clarendon, 1892, pp.68--73.
\bibitem{b3} I. S. Jacobs and C. P. Bean, ``Fine particles, thin films and exchange anisotropy,'' in Magnetism, vol. III, G. T. Rado and H. Suhl, Eds. New York: Academic, 1963, pp. 271--350.
\bibitem{b4} K. Elissa, ``Title of paper if known,'' unpublished.
\bibitem{b5} R. Nicole, ``Title of paper with only first word capitalized,'' J. Name Stand. Abbrev., in press.
\bibitem{b6} Y. Yorozu, M. Hirano, K. Oka, and Y. Tagawa, ``Electron spectroscopy studies on magneto-optical media and plastic substrate interface,'' IEEE Transl. J. Magn. Japan, vol. 2, pp. 740--741, August 1987 [Digests 9th Annual Conf. Magnetics Japan, p. 301, 1982].
\bibitem{b7} M. Young, The Technical Writer's Handbook. Mill Valley, CA: University Science, 1989.
\end{thebibliography}
\vspace{12pt}
\end{document}
